\chapter{Collections}\label{ch:collections}

A \term{collection} is an object that holds other objects, its
\term{elements}, each under a \term{key}: lists, vectors, strings,
tables, ranges and the rest. Dylan defines collections by a protocol of
generic functions rather than by their representations, so that one
function, \code{map} or \code{sort} or \code{reduce}, works on all of
them, and a class you define joins them by providing a handful of
methods. This chapter follows chapter~12 of the 1992 manual
(\dylanbook{chapter 12}) and the parts of chapter~9 that concern
collections.

The chapter first describes the classes and the protocol every
collection shares, then the functions on all collections and on
sequences, and then each built-in class in turn. Examples show their
results as comments. A list, a vector and a string print the same in
both implementations; the other collections print differently, and
where that matters the examples convert them to lists first. A class in
a result is written by its name alone, as in chapter~\ref{ch:numbers}.

\section{The collection classes}\label{sec:collections-classes}

\index{collection classes}
Collections are classified along two independent lines. A collection
either finds its elements by integer position, as a \term{sequence}
does, or by keys of any kind, as an \term{explicit-key collection} does;
and it either can be changed after it is made, as a \term{mutable
collection} can, or cannot. The abstract classes that express this, and
the concrete classes below them, are:

\begin{itemize}[nosep]
\item \code{<collection>}\index{<collection>}
  \begin{itemize}[nosep]
  \item \code{<explicit-key-collection>}\index{<explicit-key-collection>}
  \item \code{<mutable-collection>}\index{<mutable-collection>}
  \item \code{<sequence>}\index{<sequence>}
    \begin{itemize}[nosep]
    \item \code{<range>}
    \end{itemize}
  \item \code{<mutable-explicit-key-collection>}\index{<mutable-explicit-key-collection>}
    (an explicit-key and a mutable collection)
    \begin{itemize}[nosep]
    \item \code{<table>}
    \item \code{<array>}
    \end{itemize}
  \item \code{<mutable-sequence>}\index{<mutable-sequence>}
    (a mutable collection and a sequence)
    \begin{itemize}[nosep]
    \item \code{<list>}, with subclasses \code{<pair>} and \code{<empty-list>}
    \item \code{<deque>}
    \item \code{<string>}
    \item \code{<vector>} (also an \code{<array>})
      \begin{itemize}[nosep]
      \item \code{<simple-object-vector>}
      \item \code{<stretchy-vector>}
      \item \code{<byte-string>} (also a \code{<string>})
      \item \code{<unicode-string>} (also a \code{<string>})
      \end{itemize}
    \end{itemize}
  \end{itemize}
\end{itemize}

\noindent The concrete classes, the ones a program makes instances of,
are these:

\begin{center}
\begin{tabular}{p{0.32\textwidth}p{0.6\textwidth}}
\toprule
Class & What it is \\
\midrule
\code{<list>}, \code{<pair>}, \code{<empty-list>} & linked lists of pairs (sealed) \\
\code{<simple-object-vector>} & a fixed-size vector of any objects (sealed) \\
\code{<byte-string>} & a fixed-size string of 8-bit characters (sealed) \\
\code{<unicode-string>} & a fixed-size string of any characters (sealed; compiler to C only) \\
\code{<stretchy-vector>} & a vector that grows and shrinks \\
\code{<deque>} & a double-ended queue \\
\code{<range>} & an arithmetic progression, possibly unbounded; not mutable \\
\code{<table>} & a hash table from any keys to any elements \\
\code{<array>} & a multi-dimensional array \\
\bottomrule
\end{tabular}
\end{center}

\noindent \code{<vector>}, \code{<string>} and \code{<array>} are
instantiable too, but \code{make} of them answers an instance of a
concrete subclass: \code{(make <vector> size: 3)} is a
\code{<simple-object-vector>}, \code{(make <string>)} a
\code{<byte-string>}, and an \code{<array>} of one dimension is a
\code{<simple-object-vector>}.

\section{Keys, elements and size}\label{sec:collections-keys}

\index{element}\index{key}
Every collection is a mapping from keys to elements. The keys of a
sequence are the integers from 0 up to its size; the keys of a table are
whatever objects were stored in it; the keys of an array are sequences
of indices. \code{element} answers the element under a key, and
\code{(setter element)} replaces it:

\begin{dylan}
(element '(a b c) 1)                       ; => b
(element "abc" 0)                          ; => #\a
(element #(a b c) 5 default: 'none)        ; => none
(define v (vector 1 2 3))
(set! (element v 0) 'first)
v                                          ; => #(first 2 3)
\end{dylan}

\noindent Asking for a key the collection does not have is an error,
unless a \code{default:} is given, which is then the answer. Setting the
element of a key a collection does not have is an error too, except for
a table, which adds the key.

\index{size}\index{empty?}\index{key-sequence}
\code{size} answers the number of elements, or \code{\#f} for a
collection of unbounded size; \code{empty?} answers whether there are
none; and \code{key-sequence} answers a sequence of the keys:

\begin{dylan}
(size '(1 2 3))           ; => 3
(size "abc")              ; => 3
(size (range from: 1))    ; => #f
(empty? '())              ; => #t
(empty? #(1))             ; => #f
(key-sequence "ab")       ; => a sequence of 0 and 1
\end{dylan}

\noindent The key sequence of a sequence holds the integers from 0 up to
its size (a range, except that the Common Lisp implementation answers a
list for a vector or a string); that of a table is a list of its keys, in
no particular order.

\index{class-for-copy}
\code{class-for-copy} answers the class used when a function such as
\code{map} or \code{reverse} makes a new collection like an existing
one. For most collections it is the collection's own class; a range,
which is not mutable, answers \code{<list>}, and a string answers
\code{<byte-string>}. A collection class that is not mutable must define
a method for it (section~\ref{sec:collections-own}).

\section{The iteration protocol}\label{sec:collections-iteration}

\index{iteration protocol}
All iteration over collections, by \code{for-each}, \code{map},
\code{reduce} and every other function in this chapter, is done with
four generic functions. A collection chooses its own representation for
an \term{iteration state}, an object recording where an iteration has
got to; for every collection, the state \code{\#f} means that the
iteration is finished.

\begin{description}
\item[\code{(initial-state \meta{collection})}]\index{initial-state}
  answers the state for the first element, or \code{\#f} if the
  collection is empty.
\item[\code{(next-state \meta{collection} \meta{state})}]\index{next-state}
  answers the state for the next element, or \code{\#f} if there is
  none. It may change the state it is given, which must not be used
  again.
\item[\code{(current-element \meta{collection} \meta{state})}]\index{current-element}
  answers the element at a state.
\item[\code{(copy-state \meta{collection} \meta{state})}]\index{copy-state}
  answers a state that \code{next-state} will not change, for a program
  that wants to return to a place in an iteration.
\end{description}

\noindent Explicit-key collections add one more, and mutable collections
another:

\begin{description}
\item[\code{(current-key \meta{collection} \meta{state})}]\index{current-key}
  answers the key of the element at a state.
\item[\code{((setter current-element) \meta{collection} \meta{state} \meta{value})}]
  \index{current-element!setter}
  replaces the element at a state; it is written
  \code{(set! (current-element \meta{c} \meta{s}) \meta{value})}.
\end{description}

\noindent And collections that can be walked backward, the vectors, the
strings, the deques and the bounded ranges, also have:

\begin{description}
\item[\code{(final-state \meta{collection})}]\index{final-state}
  answers the state for the last element, or \code{\#f} if the
  collection is empty.
\item[\code{(previous-state \meta{collection} \meta{state})}]\index{previous-state}
  answers the state for the element before, or \code{\#f}.
\end{description}

Here is the protocol used directly. A sequence's keys are counted
alongside its states, since a sequence is walked in key order:

\begin{dylan}
(define v (vector 10 20 30))
(for ((state (initial-state v) (next-state v state))
      (key 0 (+ key 1)))
     ((not state))
  (print (list key (current-element v state))))
; prints (0 10), (1 20), (2 30)

(for ((state (final-state v) (previous-state v state)))
     ((not state))
  (print (current-element v state)))
; prints 30, 20, 10
\end{dylan}

\noindent A table supplies its keys through \code{current-key}, and any
mutable collection can be changed in place as it is walked:

\begin{dylan}
(define t (make <table>))
(set! (element t 'a) 1)
(for ((state (initial-state t) (next-state t state)))
     ((not state))
  (print (list (current-key t state) (current-element t state))))
; prints (a 1)

(define l (list 1 2 3))
(for ((state (initial-state l) (next-state l state)))
     ((not state))
  (set! (current-element l state) (* 10 (current-element l state))))
l    ; => (10 20 30)
\end{dylan}

\noindent \code{copy-state} keeps a place while the iteration goes on:

\begin{dylan}
(bind ((s (initial-state l)))
  (bind ((saved (copy-state l s)))
    (set! s (next-state l s))
    (list (current-element l s) (current-element l saved))))
; => (20 10)
\end{dylan}

\noindent The protocol assumes that a collection does not change while it
is iterated over, except through \code{(setter current-element)}. A
\term{sequence} must also be \term{stable}: every iteration visits its
elements in the same order, the order of their keys, which is its
\term{natural order}. A table need not be: two iterations over the same
table may visit its elements in different orders.

\begin{notebox}
A 1994 design note, number~26, later replaced this protocol with two
functions, \code{forward-iteration-protocol} and
\code{backward-iteration-protocol}. Both implementations follow the 1992 manual: the protocol is the one
described here, and the newer functions do not exist.
\end{notebox}

Most programs never call these functions; they use \code{for-each} and
the other loops of chapter~\ref{ch:control}, or the functions in the
rest of this chapter, all of which are built on them.

\section{Defining your own collection}\label{sec:collections-own}

\index{collection!defining}
To make a new kind of collection, define a subclass of one of the
abstract classes of section~\ref{sec:collections-classes}: a sequence
class for a collection keyed by position, an explicit-key class for one
with keys of its own. Then give it methods for the iteration protocol:

\begin{itemize}
\item every collection: \code{initial-state}, \code{next-state},
  \code{current-element} and \code{copy-state};
\item an explicit-key collection: also \code{current-key};
\item a mutable collection: also \code{(setter current-element)};
\item a collection that is not mutable: also \code{class-for-copy},
  naming a mutable class for copies.
\end{itemize}

\noindent Nothing else is required: \code{element}, \code{size},
\code{map}, \code{reduce}, \code{sort} and the rest all have default
methods that work through the protocol. They may be slow (the default
\code{element} on a sequence walks from the start), so a class that can
do better should also define methods for \code{element},
\code{(setter element)}, \code{size} or \code{member?}.

\subsection{A sequence}

This class is an immutable sequence of the first \meta{count} squares,
computed when asked for. Its iteration state is simply the index:

\begin{dylan}
(define-class <squares> (<sequence>)
  (count init-keyword: count: init-value: 0))

(define-method initial-state ((s <squares>))
  (if (> (count s) 0) 0 #f))

(define-method next-state ((s <squares>) (state <integer>))
  (bind ((next (+ state 1)))
    (if (< next (count s)) next #f)))

(define-method current-element ((s <squares>) (state <integer>))
  (* state state))

(define-method copy-state ((s <squares>) (state <integer>))
  state)

(define-method class-for-copy ((s <squares>)) <list>)
\end{dylan}

\noindent With only these, it is a full collection:

\begin{dylan}
(define sq (make <squares> count: 6))
(size sq)                          ; => 6
(element sq 4)                     ; => 16
(map (method (x) (+ x 1)) sq)      ; => (1 2 5 10 17 26)
(reduce + 0 sq)                    ; => 55
(member? 9 sq)                     ; => #t
(choose odd? sq)                   ; => (1 9 25)
(reverse sq)                       ; => (25 16 9 4 1 0)
(as <vector> sq)                   ; => #(0 1 4 9 16 25)
(= sq '(0 1 4 9 16 25))            ; => #t
(for-each ((x sq)) ()
  (display x) (display " "))       ; prints 0 1 4 9 16 25
\end{dylan}

\noindent \code{map}, \code{choose} and \code{reverse} answer lists,
because \code{class-for-copy} says so.

\subsection{A mutable explicit-key collection}

This class keeps its keys and elements in an association list, a list
of pairs. Its iteration state is the rest of that list, from the current
pair on; it supplies \code{current-key}, and, being mutable,
\code{(setter current-element)}:

\begin{dylan}
(define-class <alist> (<mutable-explicit-key-collection>)
  (entries init-value: '()))

(define-method initial-state ((a <alist>))
  (if (empty? (entries a)) #f (entries a)))
(define-method next-state ((a <alist>) (state <pair>))
  (if (empty? (cdr state)) #f (cdr state)))
(define-method current-element ((a <alist>) (state <pair>))
  (cdr (car state)))
(define-method current-key ((a <alist>) (state <pair>))
  (car (car state)))
(define-method (setter current-element) ((a <alist>) (state <pair>) value)
  (set! (cdr (car state)) value))
\end{dylan}

\noindent The default \code{element} and \code{(setter element)} find a
key by comparing with \code{=}; \code{key-sequence}, \code{size},
\code{do} and the others work as for any collection:

\begin{dylan}
(define a (make <alist>))
(set! (entries a) (list (cons 'red 1) (cons 'green 2) (cons 'blue 3)))
(element a 'green)                ; => 2
(element a 'pink default: 0)      ; => 0
(set! (element a 'blue) 30)
(element a 'blue)                 ; => 30
(size a)                          ; => 3
(reduce + 0 a)                    ; => 33
(any? (method (v) (> v 10)) a)    ; => #t
(do print a)                      ; prints 1, 2 and 30
\end{dylan}

\noindent (The class need not define \code{copy-state}: the default
method answers the state itself, which is right for a state that
\code{next-state} never changes, as here.)

\section{Functions on all collections}\label{sec:collections-functions}

\index{collection functions}
The functions in this section work on every collection. Those that take
several collections process corresponding elements together: elements
with the same key. For sequences that means elements at the same
position, and processing stops at the end of the shortest. Other
collections are \term{aligned}: only the keys all of them have are
processed (\dylanbook{chapter 12}, ``Collection Alignment'').

\subsection{do, map, map-as and map-into}

\index{do}
\code{(do \meta{function} \meta{collection} ...)} calls the function on
each group of corresponding elements, for effect, and answers
\code{\#f}:

\begin{dylan}
(do (method (a b) (print (+ a b)))
    '(100 100 200 200)
    '(1 2 3 4))
; prints 101, 102, 203 and 204, then answers #f
\end{dylan}

\index{map}\index{map-as}
\code{map} answers a new collection of the results, of the
\code{class-for-copy} of the first collection; \code{map-as} takes the
class of the result as its first argument:

\begin{dylan}
(map + '(100 100 200 200) '(1 2 3 4))              ; => (101 102 203 204)
(map + '(1 2 3) #(10 20))                          ; => (11 22)
(map + (range from: 1) '(10 20 30))                ; => (11 22 33)
(map as-uppercase "abc")                           ; => "ABC"
(map-as <vector> + '(100 100 200 200) '(1 2 3 4))  ; => #(101 102 203 204)
\end{dylan}

\noindent Mapping over a table makes a table with the same keys:

\begin{dylan}
(define prices (make <table>))
(set! (element prices 'tea) 3)
(set! (element prices 'cake) 5)
(define doubled (map (curry * 2) prices))
(element doubled 'cake)                            ; => 10
\end{dylan}

\noindent The class of the result decides what it can hold. \code{map}
over a string makes a string, so a function that answers something other
than a character is an error there; ask for another class with
\code{map-as}:

\begin{dylan}
(map-as <list> (method (c) 1) "ab")                ; => (1 1)
(map (method (c) 'x) "ab")                         ; => error
\end{dylan}

\noindent The compiler to C reports the error as \code{a <byte-string>
holds only 8-bit characters: x}; the Common Lisp implementation as
\code{as -{}- not specialized for this class type}. (The Common Lisp
implementation turns an integer result into the character with that code
instead, so \code{(map (method (c) 1) "ab")} answers a string of two
characters there and is an error in the compiler to C.)

\index{map-into}
\code{(map-into \meta{target} \meta{function} \meta{collection} ...)}
stores the results into the mutable collection \meta{target}, replacing
its elements, and answers it. The target may be one of the sources.

\begin{dylan}
(define x (list 100 100 200 200))
(map-into x + '(1 2 3 4) '(10 20 30 40))   ; => (11 22 33 44)
x                                          ; => (11 22 33 44)
(map-into x - x)                           ; => (-11 -22 -33 -44)
\end{dylan}

\noindent The function receives only the elements of the source
collections. The 1992 manual's own example,
\code{(map-into x + '(1 2 3 4))} answering \code{(101 102 203 204)},
also passes the target's elements, contrary to its prose; in both
implementations it answers \code{(1 2 3 4)} (appendix~\ref{app:differences}).
Mapping into a table stores each result under its key, adding keys the
table lacks.

\subsection{any?, every?, reduce and reduce1}

\index{any?}\index{every?}
\code{any?} answers the first true value the function returns for a
group of elements, or \code{\#f}; \code{every?} answers \code{\#t} if the
function is true for every group and \code{\#f} otherwise. Both stop as
soon as the answer is known.

\begin{dylan}
(any? > '(1 2 3 4) '(5 4 3 2))       ; => #t
(any? even? '(1 3 5 7))              ; => #f
(any? (method (x) (and (even? x) (* x 10))) '(1 2 3))   ; => 20
(every? > '(1 2 3 4) '(5 4 3 2))     ; => #f
(every? odd? '(1 3 5 7))             ; => #t
\end{dylan}

\index{reduce}\index{reduce1}
\code{(reduce \meta{function} \meta{initial} \meta{collection})}
combines the elements from the first to the last: it calls the function
on the initial value and the first element, then on that result and the
second element, and so on, and answers the last result, or the initial
value if the collection is empty. \code{reduce1} uses the first element
as the initial value; it is an error on an empty collection.

\begin{dylan}
(reduce max 10 '(3 1 4 1 5 9))    ; => 10
(reduce max 10 '(3 12 9 8 8 6))   ; => 12
(reduce - 0 '(1 2 3))             ; => -6      ((0 - 1) - 2) - 3
(reduce + 0 '())                  ; => 0
(reduce1 + '(1 2 3 4 5))          ; => 15
\end{dylan}

\subsection{member? and find-key}

\index{member?}
\code{(member? \meta{value} \meta{collection} test: \meta{test})}
answers whether the collection has an element that matches the value.
The test defaults to \code{id?} and is called with the value first and
an element second.

\begin{dylan}
(define flavors #(chocolate pistachio pumpkin))
(member? 'chocolate flavors)         ; => #t
(member? 'banana flavors)            ; => #f
(member? 2.0 #(1 2 3))               ; => #f      not id? to 2
(member? 2.0 #(1 2 3) test: =)       ; => #t
\end{dylan}

\noindent On a list, as the 1992 manual says, the true value is the rest
of the list from the matching element on: \code{(member? 2 '(1 2 3))}
answers \code{(2 3)}. (A 1993 design note made it a boolean; neither
implementation follows that.)

\index{find-key}
\code{(find-key \meta{collection} \meta{predicate})} answers the key of
the first element for which the predicate is true, or the
\code{failure:} value (default \code{\#f}). \code{skip:} \meta{n} skips
the first \meta{n} matches.

\begin{dylan}
(find-key flavors (method (f) (id? f 'pistachio)))   ; => 1
(find-key '(1 2 3 4) even?)                          ; => 1
(find-key '(1 2 3 4) even? skip: 1)                  ; => 3
(find-key '(1 2 3 4) even? skip: 2)                  ; => #f
(find-key '(1 3) even? failure: 'none)               ; => none
\end{dylan}

\subsection{replace-elements! and fill!}

\index{replace-elements"!}
\code{(replace-elements! \meta{collection} \meta{predicate}
\meta{new-value-function})} replaces each element for which the
predicate is true by the result of calling the function on it, at most
\code{count:} of them if that is given, and answers the collection:

\begin{dylan}
(replace-elements! (list 10 13 16 19) odd? (method (x) (* 2 x)))
; => (10 26 16 38)
(replace-elements! (list 1 3 5 7) odd? (method (x) 0) count: 2)
; => (0 0 5 7)
\end{dylan}

\index{fill"!}
\code{(fill! \meta{collection} \meta{value})} stores the value under
every key. For a sequence, \code{start:} (inclusive, default 0) and
\code{end:} (exclusive, default the size) limit it to a part:

\begin{dylan}
(fill! (list 'a 'b 'c 'd 'e 'f) 3 start: 2)      ; => (a b 3 3 3 3)
(fill! (vector 1 2 3 4 5) 0 start: 1 end: 3)     ; => #(1 0 0 4 5)
\end{dylan}

\section{Functions on sequences}\label{sec:collections-sequences}

\index{sequence functions}
These functions work on any sequence. Unless a description says
otherwise, a function that answers a new sequence answers one of the
\code{class-for-copy} of its (first) argument and leaves the argument
unchanged; a function whose name ends in \code{!} may change its argument
and reuse its storage, and its result may or may not be the argument
itself, so always use the result. Functions that take a \code{test:}
default it to \code{id?}.

\subsection{Adding and removing}

\index{add}\index{add"!}\index{add-new}\index{add-new"!}
\code{(add \meta{sequence} \meta{element})} answers a new sequence with
the element added; \code{add!} may add it to the sequence itself. Where
the element goes depends on the class: at the front of a list, a vector,
a string and a deque. \code{add!} on a stretchy vector adds at the end;
\code{add} on one adds at the end in the compiler to C and at the front
in the Common Lisp implementation.
\code{add-new} and \code{add-new!} add the element only if the sequence
does not already have it (by \code{test:}).

\begin{dylan}
(define numbers '(3 4 5))
(add numbers 1)             ; => (1 3 4 5)
numbers                     ; => (3 4 5)
(add #(3 4 5) 1)            ; => #(1 3 4 5)
(add "bc" #\a)              ; => "abc"
(add! (list 3 4 5) 1)       ; => (1 3 4 5)
(add-new '(3 4 5) 1)        ; => (1 3 4 5)
(add-new '(3 4 5) 4)        ; => (3 4 5)
\end{dylan}

\index{remove}\index{remove"!}\index{choose}\index{choose-by}
\code{(remove \meta{sequence} \meta{value})} answers the sequence
without the elements that match the value, or without the first
\code{count:} of them; \code{remove!} may change the sequence.
\code{choose} keeps the elements for which a predicate is true, and
\code{choose-by} keeps the elements of one sequence whose partners in
another satisfy a predicate.

\begin{dylan}
(remove '(3 1 4 1 5 9) 1)                 ; => (3 4 5 9)
(remove '(3 1 4 1 5 9) 1 count: 1)        ; => (3 4 1 5 9)
(remove "banana" #\a)                     ; => "bnn"
(remove! (list 3 1 4 1 5 9) 1)            ; => (3 4 5 9)
(choose even? '(3 1 4 1 5 9))             ; => (4)
(choose even? #(1 2 3 4))                 ; => #(2 4)
(choose-by even? (range from: 1) '(a b c d e f g h i))
; => (b d f h)
\end{dylan}

\subsection{Sets}

\index{intersection}\index{union}\index{remove-duplicates}
\index{remove-duplicates"!}
\code{intersection} answers the elements of the first sequence that are
also in the second; \code{union} answers the elements of both, those of
the second only if the first lacks them; \code{remove-duplicates}
answers a sequence without repeated elements, keeping the first of each.
The test of \code{intersection} and \code{union} is called with an
element of the first sequence and one of the second.

\begin{dylan}
(intersection '(john paul george ringo)
              '(richard george edward charles john))
; => (john george)
(union '(butter flour sugar salt eggs)
       '(eggs butter mushrooms onions salt))
; => (butter flour sugar salt eggs mushrooms onions)
(remove-duplicates '(spam eggs spam sausage spam spam spam))
; => (spam eggs sausage)
(remove-duplicates '(1 1.0 2) test: =)   ; => (1 2)
\end{dylan}

\subsection{Copying and joining}

\index{copy-sequence}\index{concatenate}\index{concatenate-as}
\index{replace-subsequence"!}
\code{copy-sequence} answers a new sequence of the elements from
\code{start:} (default 0) up to \code{end:} (default the size).
\code{concatenate} joins sequences into a new one of the first one's
\code{class-for-copy}; \code{concatenate-as} takes the class.
\code{replace-subsequence!} overwrites the elements of a sequence from
\code{start:} on with those of another, which must fit.

\begin{dylan}
(define hamlet '(to be or not to be))
(id? hamlet (copy-sequence hamlet))          ; => #f
(copy-sequence hamlet start: 2 end: 4)       ; => (or not)
(copy-sequence "hello" start: 1)             ; => "ello"
(concatenate "low-" "calorie")               ; => "low-calorie"
(concatenate '(0 1 2) '(3 4 5) '(6 7 8))     ; => (0 1 2 3 4 5 6 7 8)
(concatenate #(1) '(2))                      ; => #(1 2)
(concatenate-as <string> '(#\n #\o #\n) '(#\f #\a #\t))   ; => "nonfat"
(concatenate-as <vector> '(0 1 2) '(3 4 5))  ; => #(0 1 2 3 4 5)
(define phrase (copy-sequence "I hate oatmeal."))
(replace-subsequence! phrase "like" start: 2)  ; => "I like oatmeal."
\end{dylan}

\noindent (\code{phrase} is a copy because a string literal is a
constant; see section~\ref{sec:collections-strings}.)

\subsection{Reversing and sorting}

\index{reverse}\index{reverse"!}\index{sort}\index{sort"!}
\code{reverse} answers the elements in the opposite order and
\code{reverse!} may reuse the sequence's storage. \code{sort} answers the
elements in ascending order by \code{test:}, which defaults to \code{<};
\code{sort!} may reuse the storage. Both sorts are stable, keeping equal
elements in their original order, whether or not \code{stable:} is given.

\begin{dylan}
(define x '(bim bam boom))
(reverse x)                  ; => (boom bam bim)
x                            ; => (bim bam boom)
(reverse "abc")              ; => "cba"
(reverse! (list 1 2 3))      ; => (3 2 1)
(define nums '(3 1 4 1 5 9))
(sort nums)                  ; => (1 1 3 4 5 9)
nums                         ; => (3 1 4 1 5 9)
(sort nums test: >)          ; => (9 5 4 3 1 1)
(sort #(3 1 2))              ; => #(1 2 3)
(sort "hello")               ; => "ehllo"
(sort '("pear" "apple" "fig"))           ; => ("apple" "fig" "pear")
(sort '((b . 2) (a . 1) (c . 2) (d . 1))
      test: (method (x y) (< (cdr x) (cdr y))))
; => ((a . 1) (d . 1) (b . 2) (c . 2))
\end{dylan}

\subsection{Particular elements and searching}

\index{first}\index{second}\index{third}\index{last}
\code{first}, \code{second}, \code{third} and \code{last} answer those
elements, and with \code{set!} replace them. Asking for an element that
is not there is an error.

\begin{dylan}
(first '(a b c))              ; => a
(third '(a b c))              ; => c
(last '(emperor of china))    ; => china
(last #(1 2 3))               ; => 3
(define l (list 1 2 3))
(set! (first l) 'one)
(set! (last l) 'three)
l                             ; => (one 2 three)
\end{dylan}

\index{subsequence-position}
\code{(subsequence-position \meta{big} \meta{pattern})} answers the index
at which \meta{pattern} first occurs in \meta{big}, or the
\code{count:}th occurrence, or \code{\#f}:

\begin{dylan}
(subsequence-position "Ralph Waldo Emerson" "Waldo")   ; => 6
(subsequence-position "banana" "an" count: 2)          ; => 3
(subsequence-position "banana" "x")                    ; => #f
(subsequence-position '(a b c b c) '(b c))             ; => 1
\end{dylan}

\section{Lists and pairs}\label{sec:collections-lists}

\index{<list>}\index{<pair>}\index{<empty-list>}\index{list}
A list is either the empty list, \code{'()}, the only instance of
\code{<empty-list>}, or a \code{<pair>}, whose \term{car} is an element
and whose \term{cdr} is the rest of the list. The classes are sealed. The
empty list is not \code{\#f}: it is a true value.

\index{cons}\index{list*}\index{car}\index{cdr}
\begin{dylan}
(cons 1 2)                  ; => (1 . 2)
(cons 1 '(2 3))             ; => (1 2 3)
(list 1 2 3)                ; => (1 2 3)
(list* 1 2 3 '(4 5 6))      ; => (1 2 3 4 5 6)      the last argument is the tail
(list* 1 2)                 ; => (1 . 2)
(car '(4 5 6))              ; => 4
(cdr '(4 5 6))              ; => (5 6)
(car '())                   ; => ()
(cdr '())                   ; => ()
(object-class '())          ; => <empty-list>
(object-class '(1))         ; => <pair>
\end{dylan}

\noindent \code{(setter car)} and \code{(setter cdr)} change a pair. The
twelve functions \code{caar} through \code{cdddr} compose \code{car} and
\code{cdr} as their names spell (\code{(caddr x)} is
\code{(car (cdr (cdr x)))}), and each has a setter:

\begin{dylan}
(define x (list 4 5 6))
(set! (car x) 9)
(set! (cdr (cdr x)) '(a b))
x                          ; => (9 5 a b)
(cadr '(1 2 3))            ; => 2
(caddr '(1 2 3))           ; => 3
(cddr '(1 2 3))            ; => (3)
(define y (list 1 2 3))
(set! (cadr y) 'two)
y                          ; => (1 two 3)
\end{dylan}

\index{append}\index{find-pair}
\code{append} joins lists; its result shares the last list.
\code{(find-pair \meta{item} \meta{list})} finds the first element of a
list of pairs whose car is the item (by \code{test:}):

\begin{dylan}
(append '(1 2) '(3) '(4 5))               ; => (1 2 3 4 5)
(find-pair 'b '((a . 1) (b . 2)))         ; => (b . 2)
(find-pair 'z '((a . 1)))                 ; => #f
\end{dylan}

\noindent \code{(make <list> size: \meta{n} fill: \meta{x})} makes a list
of \meta{n} copies of \meta{x} (default \code{\#f}). \code{add!} on a list
is \code{cons} and \code{remove!} may splice the list. The \code{size}
of a circular list is \code{\#f}:

\begin{dylan}
(define circle (list 1 2 3))
(set! (cdr (cddr circle)) circle)
(size circle)    ; => #f
\end{dylan}

\begin{warnbox}
Do not let the Common Lisp implementation's read-eval-print loop print a
circular list, as it would by showing the value of the \code{set!} above
typed on its own: it prints forever. Wrap such a form in
\code{(begin ... \#f)}.
\end{warnbox}

\noindent The 1993 design note that renamed \code{cons}, \code{car} and
\code{cdr} to \code{pair}, \code{head} and \code{tail} and removed
\code{list*}, \code{append} and the \code{c...r} functions is not
followed: both implementations have the 1992 names only.

\section{Vectors and arrays}\label{sec:collections-vectors}

\index{<vector>}\index{<simple-object-vector>}\index{vector}
A vector is a one-dimensional array with constant-time access by index.
\code{\#(...)} writes a constant vector, \code{vector} makes one of its
arguments, and \code{make} makes one of a given size:

\begin{dylan}
(vector 1 2 3)                                   ; => #(1 2 3)
(make <vector> size: 3)                          ; => #(#f #f #f)
(make <simple-object-vector> size: 2 fill: 'x)   ; => #(x x)
(as <vector> '(1 2))                             ; => #(1 2)
(as <list> #(1 2))                               ; => (1 2)
\end{dylan}

\noindent A \code{<simple-object-vector>} holds any objects and cannot
change size. Its elements are read and set by \code{element}, or by the
functions of arrays below.

\index{<array>}\index{aref}\index{dimensions}
An \code{<array>} has any number of dimensions; its keys are sequences
of indices. \code{make} takes \code{dimensions:} (required) and
\code{fill:} (default \code{\#f}). \code{aref} reads the element at
indices given as separate arguments, \code{(setter aref)} sets it, and
\code{dimensions} answers the dimensions as a list. A vector is an array
of one dimension.

\begin{dylan}
(define a (make <array> dimensions: '(2 3) fill: 0))
(set! (aref a 1 2) 'x)
(aref a 1 2)            ; => x
(element a '(1 2))      ; => x
(dimensions a)          ; => (2 3)
(size a)                ; => 6
(aref #(7 8 9) 1)       ; => 8
(dimensions #(1 2 3))   ; => (3)
(define vv (vector 7 8 9))
(set! (aref vv 1) 5)
vv                      ; => #(7 5 9)
\end{dylan}

\noindent An array is iterated over in row-major order. A wrong number of
indices, or an index out of bounds, is an error. \code{map} over an
array answers an array with the same dimensions:

\begin{dylan}
(define b (map list a))
(dimensions b)          ; => (2 3)
(aref b 1 2)            ; => (x)
\end{dylan}

\section{Strings}\label{sec:collections-strings}

\index{<string>}\index{<byte-string>}\index{<unicode-string>}\index{string}
A string is a mutable sequence of characters. A \code{<byte-string>}
holds characters with codes below 256, a \code{<unicode-string>} any
characters. \code{"..."} writes a constant string; \code{make} of
\code{<string>} or \code{<byte-string>} takes \code{size:} and
\code{fill:} (default \code{\#\textbackslash{}space}).

\begin{dylan}
(make <string> size: 3 fill: #\z)   ; => "zzz"
(object-class "abc")                ; => <byte-string>
(as <list> "abc")                   ; => (#\a #\b #\c)
(as <string> '(#\h #\i))            ; => "hi"
(size "abc")                        ; => 3
\end{dylan}

\noindent The implementations differ on Unicode. In the compiler to C a
string literal with a character above code 255 is a
\code{<unicode-string>}, and a byte string cannot hold such a character
(storing one is an error). In the Common Lisp implementation every string
is a \code{<byte-string>} that holds any character, and \code{make} of
\code{<unicode-string>} is an error.

\begin{warnbox}
A string literal is a constant, and the compiler to C may share one
literal among all its occurrences, so \code{(id? "abc" "abc")} can be
\code{\#t} there (it is \code{\#f} in the Common Lisp implementation).
Do not change a literal: copy it first with \code{copy-sequence}.
\end{warnbox}

\index{as-uppercase}\index{as-lowercase}\index{as-uppercase"!}
\index{as-lowercase"!}
Strings compare lexicographically by their characters, a prefix being
less than the longer string; \code{=} compares their characters. The
case functions of characters extend to strings: \code{as-uppercase} and
\code{as-lowercase} answer new strings, and \code{as-uppercase!} and
\code{as-lowercase!} change the string they are given.

\begin{dylan}
(< "apple" "banana")       ; => #t
(< "app" "apple")          ; => #t
(< "Zebra" "apple")        ; => #t      Z (90) is less than a (97)
(= "abc" "abc")            ; => #t
(as-uppercase "Van Gogh")  ; => "VAN GOGH"
(as-lowercase "Van Gogh")  ; => "van gogh"
(define s (copy-sequence "Van Gogh"))
(as-uppercase! s)
s                          ; => "VAN GOGH"
(sort '("pear" "Apple" "fig")
      test: (method (a b) (< (as-lowercase a) (as-lowercase b))))
; => ("Apple" "fig" "pear")
\end{dylan}

\section{Stretchy vectors}\label{sec:collections-stretchy}

\index{<stretchy-vector>}
A \code{<stretchy-vector>} is a vector that grows and shrinks.
\code{add!} adds an element at its end and answers the same vector;
\code{remove!} removes elements from the vector itself. Its size can be
set: a larger size adds \code{\#f} elements and a smaller one drops
elements from the end.

\begin{dylan}
(define sv (make <stretchy-vector>))
(add! sv 'a)
(add! sv 'b)
(add! sv 'c)
(size sv)           ; => 3
(element sv 1)      ; => b
(set! (size sv) 5)
(as <list> sv)      ; => (a b c #f #f)
(set! (size sv) 2)
(as <list> sv)      ; => (a b)
(remove! sv 'a)
(as <list> sv)      ; => (b)
\end{dylan}

\noindent \code{make} takes \code{size:} and \code{fill:}. Setting an
element beyond the end is an error in both implementations; use
\code{add!}, or set the size first.

\section{Deques}\label{sec:collections-deques}

\index{<deque>}\index{push}\index{pop}\index{push-last}\index{pop-last}
A \code{<deque>}, a double-ended queue, adds and removes elements at
either end in constant time. \code{push} adds at the front and
\code{push-last} at the back, each answering the deque; \code{pop} and
\code{pop-last} remove and answer the first or the last element, and it
is an error to pop an empty deque. \code{add!} is \code{push}, and
\code{remove!} removes from the deque itself.

\begin{dylan}
(define d (make <deque>))
(push d 2)
(push d 1)
(push-last d 3)
(add! d 0)
(as <list> d)       ; => (0 1 2 3)
(pop d)             ; => 0
(pop-last d)        ; => 3
(as <list> d)       ; => (1 2)
(remove! d 2)
(as <list> d)       ; => (1)
\end{dylan}

\noindent A deque prints as \code{\{<deque> 0 1 2 3\}}.
\code{make} takes \code{size:} and \code{fill:}. The Common Lisp
implementation can also set a deque's size; the compiler to C cannot.

\section{Ranges}\label{sec:collections-ranges}

\index{<range>}\index{range}
A \code{<range>} is an arithmetic progression, computed rather than
stored, and possibly unbounded. \code{range} makes one from keywords:
\code{from:} (default 0) is the first value and \code{by:} (default 1)
the step; \code{up-to:} is an exclusive bound, \code{through:} an
inclusive one, and \code{size:} a number of elements. With none of the
last three the range is unbounded. A negative \code{by:} makes a range
that runs downward, with the bounds below it. \code{make} of
\code{<range>} takes the same keywords.

\begin{dylan}
(as <list> (range from: 1 up-to: 4))              ; => (1 2 3)
(as <list> (range from: 1 through: 4))            ; => (1 2 3 4)
(as <list> (range from: 10 through: 1 by: -3))    ; => (10 7 4 1)
(as <list> (range from: 0 by: 1/2 up-to: 2))      ; => (0 1/2 1 3/2)
(as <list> (make <range> from: 2 size: 3))        ; => (2 3 4)
(size (range from: 1))                            ; => #f
(element (range from: 10 by: 5) 100)              ; => 510
\end{dylan}

\noindent Several functions treat ranges specially so that they finish
even on an unbounded one. \code{member?} computes its answer;
\code{copy-sequence} and \code{reverse} answer ranges (an unbounded range
cannot be reversed); \code{intersection} of two ranges is a range; and
\code{=} compares two ranges by their values:

\begin{dylan}
(member? 1000000 (range))                     ; => #t
(member? 6 (range from: 1 by: 2))             ; => #f
(as <list> (copy-sequence (range from: 5) start: 2 end: 5))  ; => (7 8 9)
(as <list> (intersection (range from: 0 by: 2)
                         (range from: 0 by: 3 up-to: 20)))   ; => (0 6 12 18)
(= (range from: 1 up-to: 4) (range from: 1 through: 3))      ; => #t
(= (range from: 1 up-to: 4) '(1 2 3))                        ; => #t
\end{dylan}

\noindent A range is not mutable; its \code{class-for-copy} is
\code{<list>}, so \code{map} and \code{choose} over a range answer lists:

\begin{dylan}
(map (curry * 2) (range from: 1 through: 3))   ; => (2 4 6)
(choose odd? (range from: 1 through: 6))       ; => (1 3 5)
\end{dylan}

\noindent A range prints as \code{\{range 1 by 1 size 3\}}.

\section{Tables}\label{sec:collections-tables}

\index{<table>}\index{table}\index{remove-key"!}
A \code{<table>} maps keys of any kind to elements. Keys are compared
with \code{=} and hashed with \code{=hash}, so two strings with the same
characters are the same key, and so are \code{1} and \code{1.0}. Setting
the element of a new key adds it; \code{remove-key!} removes one (and
does nothing if the key is absent).

\begin{dylan}
(define t (make <table>))
(set! (element t 'one) 1)
(set! (element t "two") 2)
(element t 'one)                     ; => 1
(element t (copy-sequence "two"))    ; => 2
(element t 'three default: 0)        ; => 0
(size t)                             ; => 2
(remove-key! t 'one)
(key-sequence t)                     ; => ("two")
\end{dylan}

\noindent Asking for a missing key without a \code{default:} is an
error. A table is unstable under iteration: its elements come in an
order of the implementation's choosing, which may change as it grows.
A table prints as \code{\{table "two" => 2\}}.

\section{Equality: = and id?}\label{sec:collections-equality}

\index{=}\index{id?}\index{equality}
Dylan has two equalities. \code{id?} answers whether two objects are the
same object, indistinguishable by any program; \code{=} answers whether
they are equal in value, by the generic function
\code{binary=}\index{binary=}:

\begin{itemize}
\item numbers are \code{=} when they have the same mathematical value
  (chapter~\ref{ch:numbers});
\item collections are \code{=} when they have the same keys and the
  elements under each key are \code{=}; for sequences, the same elements
  in the same positions, whatever the classes of the sequences;
\item pairs, including those of dotted lists, are \code{=} when their
  cars and their cdrs are;
\item anything else is \code{=} only when it is \code{id?}.
\end{itemize}

\begin{dylan}
(= '(1 2) (list 1 2))                 ; => #t
(id? '(1 2) (list 1 2))               ; => #f
(= #(1 2) '(1 2))                     ; => #t
(= "abc" #(#\a #\b #\c))              ; => #t
(= (list 1 (vector 2)) (list 1.0 (vector 2.0)))   ; => #t
(id? 'a 'a)                           ; => #t
(id? #\a #\a)                         ; => #t
\end{dylan}

\noindent The functions that search take a \code{test:} defaulting to
\code{id?}, so pass \code{test: =} to find a value equal to another but
not the same object.

\index{=hash}
To give a class of your own a notion of equality, add a method to
\code{binary=}. If its instances are to be table keys, also add one to
\code{=hash}, which must answer the same integer for any two objects
that are \code{=}:

\begin{dylan}
(define-class <point> (<object>)
  (x init-keyword: x:)
  (y init-keyword: y:))

(define-method binary= ((a <point>) (b <point>))
  (and (= (x a) (x b)) (= (y a) (y b))))

(define-method =hash ((p <point>))
  (+ (* 31 (=hash (x p))) (=hash (y p))))

(define p1 (make <point> x: 1 y: 2))
(define p2 (make <point> x: 1 y: 2))
(id? p1 p2)                           ; => #f
(= p1 p2)                             ; => #t
(member? p2 (list p1))                ; => #f
(member? p2 (list p1) test: =)        ; => a list holding p1, a true value
(define names (make <table>))
(set! (element names p1) "home")
(element names p2)                    ; => "home"
\end{dylan}
